\documentclass[../../main.tex]{subfiles}% 注意这里的文件路径不能用 ./main.tex ,否则用latexmk编译子文件会报错
\graphicspath{{\subfix{./image/}}{\subfix{../../image/}}} % 指定图片引用路径,后续可以直接使用图片文件名
% 如果图片存放在上级目录,则使用 ../../image/ 作为备用路径

% 例如:
% \begin{figure}[H]
% \centering
% \includegraphics[width=0.3\textwidth]{图.png}
% \caption{}
% \label{figure:图}
% \end{figure}
% 注意:上述\label{}一定要放在\caption{}之后,否则引用图片序号会只会显示??.

\begin{document}

% SOURCE p=36; printed=21; object=ch01-sec02-title; region=middle
\section{\texorpdfstring{$\sigma$}{σ}-代数}\label{section:ch01-sigma-algebras}

% SOURCE p=36; printed=21; object=ch01-algebra-definition; region=middle
本节讨论作为测度定义域的集合族。

\begin{definition}[集合代数与 $\sigma$-代数]\label{definition:ch01-algebra-definition}
设 $X$ 为非空集合。$X$ 上的\textbf{集合代数}是由 $X$ 的子集构成的非空族 $\mathscr{A}$，且对有限并和补运算封闭；换言之，若 $E_1,\ldots,E_n\in\mathscr{A}$，则 $\bigcup_1^nE_j\in\mathscr{A}$；若 $E\in\mathscr{A}$，则 $E^{\mathrm{c}}\in\mathscr{A}$。对可数并封闭的代数称为 \textbf{$\sigma$-代数}。（一些作者用术语“域”和“$\sigma$-域”代替“代数”和“$\sigma$-代数”。）
\end{definition}
% SUPPLEMENT: id=v14-ch01-algebra-remark; parent=ch01-algebra-definition; kind=motivation+hypotheses; origin=independent-derivation
\begin{remark}
\textbf{动机与条件说明（重构新增）。}测度要给一批集合赋值；先要求这批集合能在基本操作下保持可用，才能讨论“或发生”“不发生”和逐步逼近。固定 $X$ 使所有补集在同一环境中计算；$\mathscr{A}\subseteq\mathscr{P}(X)$ 限定了赋值对象。$\mathscr{A}$ 非空保证可以选取一个成员，再由补和并产生 $X$ 与 $\varnothing$；若只检验正项数的有限并而允许空族，这些封闭条件会空真，无法推出这两个基本成员。

有限并封闭处理有限拼接，补封闭处理否定；交无需另列，因为 deMorgan 律把交转为并和补。$\sigma$-代数的额外条件是可数并封闭，使可数次集合操作仍停留在定义域中；这不等于任意并封闭。有限并与补不足以代替可数并：$\mathbb{N}$ 上所有有限或余有限集合构成代数，补交换两类；有限并若含余有限项，其补包含于一个有限集，否则仍有限。但是所有偶数单点集的并既无限又余无限，不在该代数中。原书规定 $X$ 非空是背景约定；在空环境上，$\{\varnothing\}$ 仍满足完全相同的代数和 $\sigma$-代数封闭规则。
\end{remark}
% END SUPPLEMENT: id=v14-ch01-algebra-remark



% SOURCE p=36; printed=21; object=ch01-algebra-intersections; region=middle-lower

\begin{remark}
注意，由于 $\bigcap_jE_j=(\bigcup_jE_j^{\mathrm{c}})^{\mathrm{c}}$，代数（分别地，$\sigma$-代数）也对有限（分别地，可数）交封闭。此外，若 $\mathscr{A}$ 为代数，则 $\varnothing\in\mathscr{A}$ 且 $X\in\mathscr{A}$，因为对 $E\in\mathscr{A}$，有 $\varnothing=E\cap E^{\mathrm{c}}$ 和 $X=E\cup E^{\mathrm{c}}$。
\end{remark}
% SUPPLEMENT: id=v14-ch01-algebra-closure-proof; parent=ch01-algebra-intersections; kind=proof; origin=independent-derivation
\begin{proof}
\textbf{封闭性断言的补写证明（重构新增）。}对非空有限指标集中的 $E_j\in\mathscr{A}$，各补集 $E_j^{\mathrm c}\in\mathscr{A}$，其有限并仍在 $\mathscr{A}$ 中，再取补仍在其中；deMorgan 律即给出各 $E_j$ 的交。若 $\mathscr{A}$ 是 $\sigma$-代数，把这一步的有限并换成可数并，就得到非空可数交封闭。由 $\mathscr{A}$ 非空取 $E\in\mathscr{A}$，先由并和补得到 $X=E\cup E^{\mathrm c}\in\mathscr{A}$，再得到 $\varnothing=X^{\mathrm c}\in\mathscr{A}$。因此空并和空交按固定环境的约定也属于代数，不需要在推导中预先假定它们已存在。
\end{proof}
% END SUPPLEMENT: id=v14-ch01-algebra-closure-proof



% SOURCE p=36; printed=21; object=ch01-disjointification; region=lower

\begin{proposition}\label{proposition:ch01-disjointification}
值得注意的是，只要代数 $\mathscr{A}$ 对可数不交并封闭，它就是 $\sigma$-代数。
\end{proposition}

\begin{proof}
事实上，设 $\{E_j\}_1^\infty\subseteq\mathscr{A}$。令
\begin{align*}
F_k=E_k\setminus\Bigl[\bigcup_1^{k-1}E_j\Bigr]
=E_k\cap\Bigl[\bigcup_1^{k-1}E_j\Bigr]^{\mathrm{c}}.
\end{align*}
则各 $F_k$ 属于 $\mathscr{A}$ 且两两不交，并有 $\bigcup_1^\infty E_j=\bigcup_1^\infty F_k$。
\end{proof}
% SUPPLEMENT: id=v14-ch01-disjointification-proof; parent=ch01-disjointification; kind=proof+motivation+hypotheses; origin=independent-derivation
\begin{proof}
\textbf{原证明细节的补全（重构新增）。}当 $k=1$ 时，前面的空并为 $\varnothing$，故 $F_1=E_1\in\mathscr{A}$。对 $k>1$，先前成员的有限并、其补及与 $E_k$ 的交都在代数中，所以 $F_k\in\mathscr{A}$。若 $j<k$，则 $F_j\subseteq E_j$，而 $F_k$ 已删除 $E_j$，故二者不交。每个 $F_k\subseteq E_k$ 给出一个方向的并集包含。反之，若 $x\in\bigcup_jE_j$，非空自然数集 $\{j:x\in E_j\}$ 有最小元 $k$；$x$ 属于 $E_k$ 且不属于此前任何 $E_j$，所以 $x\in F_k$。故两个并集相等。对可数不交并封闭的假设于是给出 $\bigcup_jE_j=\bigcup_kF_k\in\mathscr{A}$，完成 $\sigma$-代数的验证。
\end{proof}
\begin{remark}
\textbf{动机与条件说明（重构新增）。}把每个点归给第一次出现的集合，可以消除重复覆盖，使一般可数并的问题归约为不交并。代数的有限并、补、交封闭用于保证每一步删去先前有限部分后仍属于代数；可数不交并封闭只在最后一步使用。这里的可数次序保证每个位置之前只有有限项，且每个已出现的点有第一次出现的指标。原来的集合不必递增，也不必预先不交；这两条都不是隐藏假设。
\end{remark}
% END SUPPLEMENT: id=v14-ch01-disjointification-proof


\begin{remark}
这种将集合序列替换为不交序列的方法值得记住；后面将多次使用。
\end{remark}


% SOURCE p=36; printed=21; object=ch01-sigma-algebra-examples; region=bottom

\begin{examples}\label{examples:ch01-sigma-algebra-examples}
下面给出几个例子。若 $X$ 为任意集合，则 $\mathscr{P}(X)$ 和 $\{\varnothing,X\}$ 都是 $\sigma$-代数。
\end{examples}
% SUPPLEMENT: id=v14-ch01-basic-sigma-examples-proof; parent=ch01-sigma-algebra-examples; kind=proof; origin=independent-derivation
\begin{proof}
\textbf{例子的补写验证（重构新增）。}任何 $X$ 的子集的补、任何可数个子集的并仍是 $X$ 的子集，且 $\mathscr{P}(X)$ 含空集，所以幂集是 $\sigma$-代数。对 $\{\varnothing,X\}$，补交换这两个成员；可数并若有一项为 $X$ 就等于 $X$，否则各项为空而并为空。因此这也是 $\sigma$-代数。$X=\varnothing$ 时两者都退化为 $\{\varnothing\}$，验证仍成立。
\end{proof}
% END SUPPLEMENT: id=v14-ch01-basic-sigma-examples-proof


\begin{definition}[可数集或余可数集的 $\sigma$-代数]\label{definition:ch01-countable-cocountable-sigma}
若 $X$ 不可数，则
\begin{align*}
\mathscr{A}=\{E\subseteq X:E\text{ 可数或 }E^{\mathrm{c}}\text{ 可数}\}
\end{align*}
是一个 $\sigma$-代数，称为\textbf{可数集或余可数集的 $\sigma$-代数}。
\end{definition}

\begin{remark}
（这里的关键在于，若 $\{E_j\}_1^\infty\subseteq\mathscr{A}$，则当所有 $E_j$ 都可数时，$\bigcup_1^\infty E_j$ 可数，否则它是余可数的。）
\end{remark}
% SUPPLEMENT: id=v14-ch01-countable-cocountable-proof; parent=ch01-countable-cocountable-sigma; kind=proof+motivation+hypotheses; origin=independent-derivation
\begin{proof}
\textbf{补写验证（重构新增）。}空集可数，故 $\varnothing\in\mathscr{A}$；补集运算交换“本身可数”和“补集可数”，故对补封闭。给定 $E_n\in\mathscr{A}$，若每项都可数，可数个可数集之并仍可数，故 $\bigcup_nE_n\in\mathscr{A}$。否则，某个 $E_k$ 不可数，按成员条件 $E_k^{\mathrm c}$ 必可数；而 $(\bigcup_nE_n)^{\mathrm c}\subseteq E_k^{\mathrm c}$，可数集的子集仍可数，所以并集余可数，仍属于 $\mathscr{A}$。这核验了 $\sigma$-代数所需全部条件。
\end{proof}
\begin{remark}
\textbf{动机与条件说明（重构新增）。}这个例子把可数部分视为“小”，把它的补视为“大”；同时纳入两类才能对补封闭。可数并封闭依赖“可数个可数集的并仍可数”和补集对包含关系的反向性。$X$ 不可数使这两类具有非平凡的区分；它并不是 $\sigma$-代数结论本身的必要条件，因为 $X$ 可数时每个子集都可数，所给集合族正好等于 $\mathscr{P}(X)$。在不可数环境中，一个集合不能同时可数且余可数，否则它与补的并 $X$ 也可数，与假设矛盾。
\end{remark}
% END SUPPLEMENT: id=v14-ch01-countable-cocountable-proof



% SOURCE p=37; printed=22; object=ch01-generated-sigma-algebra; region=top
\begin{proposition}\label{proposition:ch01-intersection-sigma-algebras}
容易验证，$X$ 上任意一族 $\sigma$-代数的交仍为 $\sigma$-代数。
\end{proposition}
% SUPPLEMENT: id=v14-ch01-intersection-sigma-proof; parent=ch01-generated-sigma-algebra; kind=proof+motivation+hypotheses; origin=independent-derivation
\begin{proof}
\textbf{补写证明（重构新增）。}设这些 $\sigma$-代数为 $\{\mathscr{M}_\alpha\}_{\alpha\in A}$，先设 $A\ne\varnothing$，记其交为 $\mathscr{D}$。每个 $\mathscr{M}_\alpha$ 都含 $X$，故 $X\in\mathscr{D}$。若 $E\in\mathscr{D}$，则 $E$ 在每个 $\mathscr{M}_\alpha$ 中，各自对补封闭，故 $E^{\mathrm c}$ 也在每个成员中，即在 $\mathscr{D}$ 中。若 $E_n\in\mathscr{D}$，对任意固定 $\alpha$，所有 $E_n$ 都在 $\mathscr{M}_\alpha$ 中，故其可数并也在其中；因 $\alpha$ 任意，该并属于 $\mathscr{D}$。这证明了结论。若 $A=\varnothing$，在环境 $\mathscr{P}(X)$ 内取空交，得到 $\mathscr{P}(X)$，它也是 $\sigma$-代数。
\end{proof}
\begin{remark}
\textbf{动机与条件说明（重构新增）。}交保留所有给定集合族同时允许的集合；封闭性条件是“对其中任意成员，运算结果仍在其中”，因此可以在每个集合族内分别检验。共同环境 $X$ 保证各处补集运算一致；指标族无需可数，因为此处是在集合族之间取交，并没有要求每个 $\sigma$-代数自身对不可数个集合的交封闭。这个区别正是构造最小生成 $\sigma$-代数的关键。
\end{remark}
% END SUPPLEMENT: id=v14-ch01-intersection-sigma-proof

\begin{definition}[生成的 $\sigma$-代数]\label{definition:ch01-generated-sigma-algebra}
因此，若 $\mathscr{E}$ 是 $\mathscr{P}(X)$ 的任意子集，则存在唯一最小的、包含 $\mathscr{E}$ 的 $\sigma$-代数 $\mathscr{M}(\mathscr{E})$，即所有包含 $\mathscr{E}$ 的 $\sigma$-代数之交。（至少总有一个这样的 $\sigma$-代数，即 $\mathscr{P}(X)$。）称 $\mathscr{M}(\mathscr{E})$ 为由 $\mathscr{E}$ \textbf{生成}的 $\sigma$-代数。
\end{definition}
% SUPPLEMENT: id=v14-ch01-generated-sigma-proof; parent=ch01-generated-sigma-algebra; kind=proof+motivation+hypotheses; origin=independent-derivation
\begin{proof}
\textbf{存在性与唯一性的补写证明（重构新增）。}所有包含 $\mathscr{E}$ 的 $\sigma$-代数形成 $\mathscr{P}(\mathscr{P}(X))$ 的子集；它包含 $\mathscr{P}(X)$，故非空。由刚证明的交封闭结论，它们的交 $\mathscr{D}$ 是 $\sigma$-代数。每个被取交的集合族都含 $\mathscr{E}$，所以 $\mathscr{E}\subseteq\mathscr{D}$。若 $\mathscr{N}$ 是任意包含 $\mathscr{E}$ 的 $\sigma$-代数，它参与这个交，故 $\mathscr{D}\subseteq\mathscr{N}$，证明最小性。若 $\mathscr{D}'$ 也具有同样最小性，分别把对方当作包含 $\mathscr{E}$ 的 $\sigma$-代数，即得 $\mathscr{D}\subseteq\mathscr{D}'$ 和 $\mathscr{D}'\subseteq\mathscr{D}$，所以二者相等。
\end{proof}
\begin{remark}
\textbf{动机与条件说明（重构新增）。}生成过程只强制加入从指定基本集合出发必需的可测集合，不任意扩大到整个幂集。唯一的输入条件是 $\mathscr{E}\subseteq\mathscr{P}(X)$；生成族本身无需非空，也无需已对任何运算封闭。$\mathscr{E}=\varnothing$ 时，所有 $\sigma$-代数都含 $\varnothing,X$，而这两个集合已构成一个 $\sigma$-代数，故生成结果为 $\{\varnothing,X\}$。这里“最小”按包含比较，强于仅仅不存在更小候选的“极小”；上面的交构造直接保证被每个候选包含。
\end{remark}
% END SUPPLEMENT: id=v14-ch01-generated-sigma-proof

下述观察常常有用。

% SOURCE p=37; printed=22; object=ch01-generated-inclusion; number=1.1; region=upper
\begin{lemma}\label{lemma:ch01-generated-inclusion}
若 $\mathscr{E}\subseteq\mathscr{M}(\mathscr{F})$，则 $\mathscr{M}(\mathscr{E})\subseteq\mathscr{M}(\mathscr{F})$。
\end{lemma}

% SOURCE p=37; printed=22; object=ch01-proof-generated-inclusion; region=upper
\begin{proof}
$\mathscr{M}(\mathscr{F})$ 是包含 $\mathscr{E}$ 的 $\sigma$-代数，因此包含 $\mathscr{M}(\mathscr{E})$。
\end{proof}
% SUPPLEMENT: id=v14-ch01-generated-inclusion-remark; parent=ch01-generated-inclusion; kind=motivation+hypotheses; origin=independent-derivation
\begin{remark}
\textbf{动机与条件说明（重构新增）。}要比较两个生成的 $\sigma$-代数，只需检查较小一边的生成元是否已在另一边中，无须列出所有反复运算的结果。关键条件 $\mathscr{E}\subseteq\mathscr{M}(\mathscr{F})$ 正好使后者成为定义 $\mathscr{M}(\mathscr{E})$ 时的一个候选，证明于是直接使用最小性。不要求 $\mathscr{E}\subseteq\mathscr{F}$，因为生成元可能由 $\mathscr{F}$ 经过运算获得。这个条件也由结论蕴含：$\mathscr{E}\subseteq\mathscr{M}(\mathscr{E})\subseteq\mathscr{M}(\mathscr{F})$。两族仍须是同一环境 $X$ 的子集族，否则这个候选比较需要先说明环境。
\end{remark}
% END SUPPLEMENT: id=v14-ch01-generated-inclusion-remark


% SOURCE p=37; printed=22; object=ch01-borel-definition; region=upper-middle

\begin{definition}[Borel $\sigma$-代数与 Borel 集]\label{definition:ch01-borel-definition}
若 $X$ 是任意度量空间，或更一般地是任意拓扑空间（见\cref{chapter:ch04-topology}），则由 $X$ 中的开集族（等价地，由 $X$ 中的闭集族）生成的 $\sigma$-代数称为 $X$ 上的 \textbf{Borel $\sigma$-代数}，记为 $\mathscr{B}_X$。它的成员称为 \textbf{Borel 集}。
\end{definition}
% SUPPLEMENT: id=v14-ch01-borel-definition-proof; parent=ch01-borel-definition; kind=proof+motivation+hypotheses; origin=independent-derivation
\begin{proof}
\textbf{开集与闭集生成等价的补写证明（重构新增）。}记开集族为 $\mathscr{O}$，闭集族为 $\mathscr{C}$。每个闭集都是某个开集相对于 $X$ 的补，所以属于 $\mathscr{M}(\mathscr{O})$，最小性给出 $\mathscr{M}(\mathscr{C})\subseteq\mathscr{M}(\mathscr{O})$。反过来，每个开集都是闭集的补，同样得到 $\mathscr{M}(\mathscr{O})\subseteq\mathscr{M}(\mathscr{C})$。两边相等。这只使用拓扑中开闭集互为补和 $\sigma$-代数的补封闭。
\end{proof}
\begin{remark}
\textbf{动机与条件说明（重构新增）。}Borel 集把空间的开集所表达的局部信息扩展为可以反复做可数集合运算的对象。需要指定拓扑，因为同一个底集配不同拓扑可能产生不同 Borel 集；具体度量只是产生拓扑的一种方式，不是定义必须保留的额外数据。这里不要求可分、完备或有限维。得到的是包含所有开集的最小 $\sigma$-代数，不保证是所有子集，也不自动包含任意测度下所有零测集的子集；后者涉及另外的完备化要求。
\end{remark}
% END SUPPLEMENT: id=v14-ch01-borel-definition-proof

因此，$\mathscr{B}_X$ 包含开集、闭集、开集的可数交、闭集的可数并，等等。

% SOURCE p=37; printed=22; object=ch01-borel-hierarchy; region=middle
这一层级的各层有标准术语。
\begin{definition}[Borel 集层级的记号]\label{definition:ch01-borel-hierarchy}
开集的可数交称为 $G_\delta$ 集；闭集的可数并称为 $F_\sigma$ 集；$G_\delta$ 集的可数并称为 $G_{\delta\sigma}$ 集；$F_\sigma$ 集的可数交称为 $F_{\sigma\delta}$ 集；如此等等。
\end{definition}
% SUPPLEMENT: id=v14-ch01-borel-hierarchy-remark; parent=ch01-borel-hierarchy; kind=motivation+hypotheses; origin=independent-derivation
\begin{remark}
\textbf{动机与条件说明（重构新增）。}这些记号记录由开闭集出发进行了哪一种可数运算，便于描述极限过程所产生的集合。$G_\delta$ 使用开集与可数交；$F_\sigma$ 使用闭集与可数并；再添的 $\sigma$ 或 $\delta$ 分别表示接着取可数并或交。每一步都在 $\sigma$-代数允许的范围内，因此这些类的所有成员都是 Borel 集。可数包括有限的情形：有限并可补入空集，有限交可补入 $X$；$\varnothing,X$ 在任何拓扑中都既开又闭。这些名称描述一种表示方式，不声称各类互不相交；例如开且闭的集合同时属于前两类。运算次序不能任意对换，定义本身没有作出这种交换断言。
\end{remark}
% END SUPPLEMENT: id=v14-ch01-borel-hierarchy-remark

（$\delta$ 和 $\sigma$ 分别代表德语 \emph{Durchschnitt} 和 \emph{Summe}，即交与并。）

% SOURCE p=37; printed=22; object=ch01-real-borel-generators-intro; region=middle
$\mathbb{R}$ 上的 Borel $\sigma$-代数在下文中将起基础性作用。为以后引用方便，我们指出它可以用多种不同的方式生成。

% SOURCE p=37; printed=22; object=ch01-real-borel-generators; number=1.2; region=middle
\begin{proposition}\label{proposition:ch01-real-borel-generators}
下列每个集合族都生成 $\mathscr{B}_{\mathbb{R}}$：
\begin{enumerate}[label=\alph*.,ref=(\alph*)]
\item\label{item:ch01-borel-open-intervals} 开区间：$\mathscr{E}_1=\{(a,b):a<b\}$；
\item\label{item:ch01-borel-closed-intervals} 闭区间：$\mathscr{E}_2=\{[a,b]:a<b\}$；
\item\label{item:ch01-borel-half-open-intervals} 半开区间：$\mathscr{E}_3=\{(a,b]:a<b\}$ 或 $\mathscr{E}_4=\{[a,b):a<b\}$；
\item\label{item:ch01-borel-open-rays} 开射线：$\mathscr{E}_5=\{(a,\infty):a\in\mathbb{R}\}$ 或 $\mathscr{E}_6=\{(-\infty,a):a\in\mathbb{R}\}$；
\item\label{item:ch01-borel-closed-rays} 闭射线：$\mathscr{E}_7=\{[a,\infty):a\in\mathbb{R}\}$ 或 $\mathscr{E}_8=\{(-\infty,a]:a\in\mathbb{R}\}$。
\end{enumerate}
\end{proposition}

% SOURCE p=37; printed=22; object=ch01-proof-real-borel-generators; region=lower
\begin{proof}
当 $j\ne3,4$ 时，$\mathscr{E}_j$ 的元素是开集或闭集；而 $\mathscr{E}_3$ 和 $\mathscr{E}_4$ 的元素是 $G_\delta$ 集，例如 $(a,b]=\bigcap_1^\infty(a,b+n^{-1})$。所有这些都是 Borel 集，所以由\cref{lemma:ch01-generated-inclusion}，对所有 $j$ 都有 $\mathscr{M}(\mathscr{E}_j)\subseteq\mathscr{B}_{\mathbb{R}}$。另一方面，$\mathbb{R}$ 中每个开集都是开区间的可数并，故再由\cref{lemma:ch01-generated-inclusion}，$\mathscr{B}_{\mathbb{R}}\subseteq\mathscr{M}(\mathscr{E}_1)$。对于 $j\geqslant2$，只需证明所有开区间都属于 $\mathscr{M}(\mathscr{E}_j)$，再应用\cref{lemma:ch01-generated-inclusion}，即可得 $\mathscr{B}_{\mathbb{R}}\subseteq\mathscr{M}(\mathscr{E}_j)$。例如，$(a,b)=\bigcup_1^\infty[a+n^{-1},b-n^{-1}]\in\mathscr{M}(\mathscr{E}_2)$。其余情形的验证留给读者（\cref{exercise:ch01-borel-generators-proof}）。
\end{proof}
% SUPPLEMENT: id=v14-ch01-real-borel-generators-remark; parent=ch01-real-borel-generators; kind=motivation+hypotheses; origin=independent-derivation
\begin{remark}
\textbf{动机与条件说明（重构新增）。}测度和可测函数的检验常从区间或单侧阈值开始，本结果说明这些简单测试已经生成实线全部 Borel 信息。此处环境是带通常拓扑的 $\mathbb{R}$；$a,b$ 为有限实数且 $a<b$，确保所列区间确为指定类型的非退化生成元。射线只需一个实端点。各结论都是分别生成同一个 $\sigma$-代数，不要求把八族一起作为生成元。证明利用可数运算把开闭端点互换，以及实线开集的可数区间表示；不能仅由一个包含方向推出生成相等。

\textbf{证明补全定位（重构新增）。}\hyperref[solution:ch01-borel-generators-proof]{对应习题的完整解答}已经逐一给出八族的两向包含和端点验证，其中闭区间逼近选取始终满足左端严格小于右端的区间，避免原简式中小指标可能出现的空区间或退化区间解释问题。该解答只使用上述定义、集合运算和有理数稠密性，不以本命题的结论证明本命题。
\end{remark}
% END SUPPLEMENT: id=v14-ch01-real-borel-generators-remark


% SOURCE p=37; printed=22; object=ch01-product-sigma-algebra-start; region=bottom

\begin{definition}[乘积 $\sigma$-代数]\label{definition:ch01-product-sigma-algebra-start}
设 $\{X_\alpha\}_{\alpha\in A}$ 为一族带指标的非空集合，$X=\prod_{\alpha\in A}X_\alpha$，且 $\pi_\alpha:X\to X_\alpha$ 为坐标映射。若对每个 $\alpha$，$\mathscr{M}_\alpha$ 都是 $X_\alpha$ 上的 $\sigma$-代数，则 $X$ 上的\textbf{乘积 $\sigma$-代数}是由
\begin{align*}
\{\pi_\alpha^{-1}(E_\alpha):E_\alpha\in\mathscr{M}_\alpha,\ \alpha\in A\}
\end{align*}
生成的 $\sigma$-代数。记此 $\sigma$-代数为 $\bigotimes_{\alpha\in A}\mathscr{M}_\alpha$。（若 $A=\{1,\ldots,n\}$，也记为 $\bigotimes_1^n\mathscr{M}_j$ 或 $\mathscr{M}_1\otimes\cdots\otimes\mathscr{M}_n$。）
\end{definition}
% SUPPLEMENT: id=v14-ch01-product-sigma-remark; parent=ch01-product-sigma-algebra-start; kind=motivation+hypotheses; origin=independent-derivation
\begin{remark}
\textbf{动机与条件说明（重构新增）。}一个柱集 $\pi_\alpha^{-1}(E_\alpha)$ 只检查第 $\alpha$ 个坐标是否落入指定集合，其余坐标不受限制。将所有这种单坐标检验生成一个 $\sigma$-代数，便得到能同时承载各坐标可测信息的最小集合族。各 $\mathscr{M}_\alpha$ 必须位于相应的 $X_\alpha$ 上，才使原像的类型匹配；坐标值的可选性来自各 $X_\alpha$ 非空，任意指标族的乘积非空使用本书的选择公理。此定义不要求指标集可数；可数性只会在把无限多个坐标限制合并为一个交时产生额外作用。若 $A=\varnothing$，乘积为只含空函数的单元素集，没有柱集生成元，生成的 $\sigma$-代数就是这个单元素空间上的 $\{\varnothing,X\}$。
\end{remark}
% END SUPPLEMENT: id=v14-ch01-product-sigma-remark

此定义的意义将在\cref{section:ch02-measurable-functions}中更为清楚；
% SOURCE p=38; printed=23; object=ch01-product-sigma-algebra-end; region=top
目前，我们先在因子个数可数的情形下，给出乘积 $\sigma$-代数的另一种、也许更直观的刻画。

% SOURCE p=38; printed=23; object=ch01-countable-product-generators; number=1.3; region=upper
\begin{proposition}\label{proposition:ch01-countable-product-generators}
若 $A$ 可数，则 $\bigotimes_{\alpha\in A}\mathscr{M}_\alpha$ 是由 $\{\prod_{\alpha\in A}E_\alpha:E_\alpha\in\mathscr{M}_\alpha\}$ 生成的 $\sigma$-代数。
\end{proposition}

% SOURCE p=38; printed=23; object=ch01-proof-countable-product-generators; region=upper
\begin{proof}
若 $E_\alpha\in\mathscr{M}_\alpha$，则 $\pi_\alpha^{-1}(E_\alpha)=\prod_{\beta\in A}E_\beta$，其中当 $\beta\ne\alpha$ 时 $E_\beta=X_\beta$；另一方面，$\prod_{\alpha\in A}E_\alpha=\bigcap_{\alpha\in A}\pi_\alpha^{-1}(E_\alpha)$。因此，结论由\cref{lemma:ch01-generated-inclusion}得到。
\end{proof}
% SUPPLEMENT: id=v14-ch01-countable-product-proof; parent=ch01-countable-product-generators; kind=proof+motivation+hypotheses; origin=independent-derivation
\begin{proof}
\textbf{两个包含方向的补全（重构新增）。}记柱集生成的 $\sigma$-代数为 $\mathscr{P}$，题设矩形族生成的为 $\mathscr{Q}$。每个柱集都能写成一个矩形：指定坐标取 $E_\alpha$，其他坐标取整个 $X_\beta\in\mathscr{M}_\beta$。所以 $\mathscr{Q}$ 含全部柱集，最小性给出 $\mathscr{P}\subseteq\mathscr{Q}$。反之，若所有 $E_\alpha\in\mathscr{M}_\alpha$，一个点属于矩形当且仅当它的每个坐标满足对应限制，故矩形为所有 $\pi_\alpha^{-1}(E_\alpha)$ 的交。$A$ 可数使该交属于 $\mathscr{P}$，所以 $\mathscr{Q}\subseteq\mathscr{P}$。若 $A$ 为空，唯一的空指标矩形为 $X$，它生成的同样是 $\{\varnothing,X\}$。
\end{proof}
\begin{remark}
\textbf{动机与条件说明（重构新增）。}单坐标限制和同时限制所有坐标给出了两种自然描述。可数性只在反向包含中把整个矩形写成可数交时使用；“每个因子是 $\sigma$-代数”保证可用全集填补其他坐标。不可数指标下不能直接把这个证明中的可数交换成不可数交。

确有失败例。设 $A$ 不可数，每个 $X_\alpha=\{0,1\}$ 且 $\mathscr{M}_\alpha=\mathscr{P}(X_\alpha)$。考虑所有成员资格只依赖于可数多个坐标的集合组成的族 $\mathscr{D}$：具体说，对其中一个集合 $E$，存在可数 $J\subseteq A$，使两个点若在 $J$ 上坐标相同，则同时属于或同时不属于 $E$。全集对应 $J=\varnothing$；补集沿用同一 $J$；可数并使用各 $J$ 的可数并，故 $\mathscr{D}$ 是 $\sigma$-代数。所有柱集属于它，所以乘积 $\sigma$-代数包含于 $\mathscr{D}$。但矩形 $\prod_\alpha\{0\}$ 不在 $\mathscr{D}$：对任意可数 $J$，取 $\beta\in A\setminus J$，全零点与只在 $\beta$ 坐标取 $1$ 的点在 $J$ 上相同，只有前者属于该矩形。故此矩形不属于乘积 $\sigma$-代数，验证了任意不可数指标推广一般不成立。
\end{remark}
% END SUPPLEMENT: id=v14-ch01-countable-product-proof


% SOURCE p=38; printed=23; object=ch01-product-generating-families; number=1.4; region=upper-middle
\begin{proposition}\label{proposition:ch01-product-generating-families}
设 $\mathscr{M}_\alpha$ 由 $\mathscr{E}_\alpha$ 生成，其中 $\alpha\in A$。则 $\bigotimes_{\alpha\in A}\mathscr{M}_\alpha$ 由 $\mathscr{F}_1=\{\pi_\alpha^{-1}(E_\alpha):E_\alpha\in\mathscr{E}_\alpha,\ \alpha\in A\}$ 生成。若 $A$ 可数，且对所有 $\alpha$ 都有 $X_\alpha\in\mathscr{E}_\alpha$，则 $\bigotimes_{\alpha\in A}\mathscr{M}_\alpha$ 由 $\mathscr{F}_2=\{\prod_{\alpha\in A}E_\alpha:E_\alpha\in\mathscr{E}_\alpha\}$ 生成。
\end{proposition}

% SOURCE p=38; printed=23; object=ch01-proof-product-generating-families; region=middle
\begin{proof}
显然，$\mathscr{M}(\mathscr{F}_1)\subseteq\bigotimes_{\alpha\in A}\mathscr{M}_\alpha$。另一方面，对每个 $\alpha$，容易看出族 $\{E\subseteq X_\alpha:\pi_\alpha^{-1}(E)\in\mathscr{M}(\mathscr{F}_1)\}$ 是 $X_\alpha$ 上包含 $\mathscr{E}_\alpha$ 的 $\sigma$-代数，从而包含 $\mathscr{M}_\alpha$。换言之，对所有 $E\in\mathscr{M}_\alpha$、$\alpha\in A$，都有 $\pi_\alpha^{-1}(E)\in\mathscr{M}(\mathscr{F}_1)$，因而 $\bigotimes_{\alpha\in A}\mathscr{M}_\alpha\subseteq\mathscr{M}(\mathscr{F}_1)$。仿照\cref{proposition:ch01-countable-product-generators}的证明，可由第一个断言推出第二个断言。
\end{proof}
% SUPPLEMENT: id=v14-ch01-product-generating-proof; parent=ch01-product-generating-families; kind=proof+motivation+hypotheses; origin=independent-derivation
\begin{proof}
\textbf{省略验证的补全（重构新增）。}置 $\mathscr{N}=\mathscr{M}(\mathscr{F}_1)$，固定 $\alpha$，记原证明中的集合族为 $\mathscr{D}_\alpha$。因为 $\pi_\alpha^{-1}(X_\alpha)=X\in\mathscr{N}$，它含 $X_\alpha$。原像保持补和可数并，故若 $E\in\mathscr{D}_\alpha$，则 $\pi_\alpha^{-1}(X_\alpha\setminus E)=X\setminus\pi_\alpha^{-1}(E)\in\mathscr{N}$；若 $E_n\in\mathscr{D}_\alpha$，则 $\pi_\alpha^{-1}(\bigcup_nE_n)=\bigcup_n\pi_\alpha^{-1}(E_n)\in\mathscr{N}$。所以它确为 $\sigma$-代数。它含 $\mathscr{E}_\alpha$，因为这些集合的原像正是 $\mathscr{F}_1$ 的生成元；于是它含 $\mathscr{M}_\alpha=\mathscr{M}(\mathscr{E}_\alpha)$。这给出所有完整柱集都在 $\mathscr{N}$ 中，最小性完成原文的反向包含。

对第二个断言，置 $\mathscr{Q}=\mathscr{M}(\mathscr{F}_2)$。每个 $\mathscr{F}_1$ 柱集都能用其他坐标填入 $X_\beta\in\mathscr{E}_\beta$ 写成 $\mathscr{F}_2$ 的矩形，故 $\mathscr{M}(\mathscr{F}_1)\subseteq\mathscr{Q}$。反之，每个 $\mathscr{F}_2$ 矩形是相应 $\mathscr{F}_1$ 柱集的可数交，因 $A$ 可数，属于 $\mathscr{M}(\mathscr{F}_1)$。于是两者相等，再用已证第一个断言得到所需乘积 $\sigma$-代数。空指标时，两种生成方式都给出单元素乘积空间的 $\{\varnothing,X\}$。
\end{proof}
\begin{remark}
\textbf{动机与条件说明（重构新增）。}无需使用每个因子中所有可测集，只使用各自的生成元已经足够；原因是原像与 $\sigma$-代数运算相容。第一个断言对指标数没有限制，也不要求生成族含全集。第二个断言另外使用可数性来取交，使用 $X_\alpha\in\mathscr{E}_\alpha$ 来把柱集填成规定形式的矩形。

第二个条件一般不能直接删去。例如两个因子均为 $\{0,1\}$，各取生成族 $\mathscr{E}_j=\{\{0\}\}$，它生成整个二点幂集。但 $\mathscr{F}_2$ 只含单点矩形 $\{(0,0)\}$，生成的集合族只有空集、该单点、其补和全集四个成员；柱集 $\{0\}\times\{0,1\}$ 不在这四者中，却属于乘积 $\sigma$-代数。因此仅有“各族分别生成因子”不足以保证所给矩形族生成乘积。加入各因子全集不会改变其生成的 $\sigma$-代数，却使填充柱集的步骤合法。
\end{remark}
% END SUPPLEMENT: id=v14-ch01-product-generating-proof


% SOURCE p=38; printed=23; object=ch01-product-borel; number=1.5; region=middle
\begin{proposition}\label{proposition:ch01-product-borel}
设 $X_1,\ldots,X_n$ 为度量空间，$X=\prod_1^nX_j$ 配备乘积度量。则 $\bigotimes_1^n\mathscr{B}_{X_j}\subseteq\mathscr{B}_X$。若各 $X_j$ 可分，则 $\bigotimes_1^n\mathscr{B}_{X_j}=\mathscr{B}_X$。
\end{proposition}

% SOURCE p=38; printed=23; object=ch01-proof-product-borel; region=middle-lower
\begin{proof}
由\cref{proposition:ch01-product-generating-families}，$\bigotimes_1^n\mathscr{B}_{X_j}$ 由集合 $\pi_j^{-1}(U_j)$（$1\leqslant j\leqslant n$）生成，其中 $U_j$ 在 $X_j$ 中为开集。由于这些集合在 $X$ 中为开集，\cref{lemma:ch01-generated-inclusion}蕴含 $\bigotimes_1^n\mathscr{B}_{X_j}\subseteq\mathscr{B}_X$。现设 $C_j$ 是 $X_j$ 中的可数稠密集，令 $\mathscr{E}_j$ 为 $X_j$ 中以 $C_j$ 中的点为球心、以有理数为半径的所有球组成的族。于是，$X_j$ 中每个开集都是 $\mathscr{E}_j$ 中成员的并；事实上是可数并，因为 $\mathscr{E}_j$ 本身可数。此外，$X$ 中对所有 $j$ 第 $j$ 个坐标均属于 $C_j$ 的点构成 $X$ 的可数稠密子集，而 $X$ 中半径为 $r$ 的球恰是各 $X_j$ 中半径为 $r$ 的球的乘积。因此，$\mathscr{B}_{X_j}$ 由 $\mathscr{E}_j$ 生成，$\mathscr{B}_X$ 由 $\{\prod_1^nE_j:E_j\in\mathscr{E}_j\}$ 生成。故由\cref{proposition:ch01-product-generating-families}，$\mathscr{B}_X=\bigotimes_1^n\mathscr{B}_{X_j}$。
\end{proof}
% SUPPLEMENT: id=v14-ch01-product-borel-proof; parent=ch01-product-borel; kind=proof+motivation+hypotheses; origin=independent-derivation
\begin{proof}
\textbf{开集基步骤的补全（重构新增）。}本书有限乘积度量为各坐标距离的最大值。因而若 $\pi_j(x)\in U_j$ 且 $U_j$ 开，可取 $r>0$ 使第 $j$ 坐标的半径 $r$ 球包含于 $U_j$；乘积中以 $x$ 为中心的半径 $r$ 球于是包含于 $\pi_j^{-1}(U_j)$。这证明柱集开，完成原文第一个包含所需的开性验证。

在可分情形，每个因子选可数稠密集 $C_j$，取正有理半径的球族 $\mathscr{E}_j$。它可数，且是开集基：若 $x_j\in U_j$ 且 $B(\varepsilon,x_j)\subseteq U_j$，稠密性给出 $c\in C_j$ 满足 $\rho_j(c,x_j)<\varepsilon/3$；再取正有理数 $r$ 使 $\rho_j(c,x_j)<r<\varepsilon-\rho_j(c,x_j)$。则 $x_j\in B(r,c)$，并由三角不等式得 $B(r,c)\subseteq B(\varepsilon,x_j)\subseteq U_j$。

所有矩形 $\prod_{j=1}^nV_j$，其中 $V_j\in\mathscr{E}_j$，组成可数族，因为只有有限个可数因子。每个矩形是开集，也是各 $\pi_j^{-1}(V_j)$ 的有限交，故在乘积 $\sigma$-代数中。它们构成 $X$ 的开集基：给定 $x\in U$ 且 $B(r,x)\subseteq U$，在每个因子中按上一步选 $x_j\in V_j\subseteq B(r,x_j)$，则 $x\in\prod_jV_j\subseteq B(r,x)\subseteq U$。所以 $U$ 是这个可数族中所有包含于 $U$ 的矩形的并，属于乘积 $\sigma$-代数。所有开集都属于其中，便得到 $\mathscr{B}_X\subseteq\bigotimes_j\mathscr{B}_{X_j}$，与原来的包含合并即为相等。这个基论证无需假定整个无界空间本身是某个有限半径球。
\end{proof}
\begin{remark}
\textbf{动机与条件说明（重构新增）。}一边由坐标可测信息生成，另一边由乘积空间开集生成；命题说明它们何时相同。有限个因子使最大值定义乘积度量，且有限个基球的矩形仍形成可数基。可分性在本证明中提供可数稠密集，从而把开集表达为可数并；它是这里使用的充分条件，不能仅因这个基证明不适用于不可分情形，就宣称每个不可分实例都不相等。第一包含只用坐标投影的开原像，不使用可分性。度量的具体最大值形式使球等于同半径坐标球的乘积；换用产生同一有限乘积拓扑的度量，Borel 集族不变。
\end{remark}
% END SUPPLEMENT: id=v14-ch01-product-borel-proof


% SOURCE p=38; printed=23; object=ch01-euclidean-product-borel; number=1.6; region=lower
\begin{corollary}\label{corollary:ch01-euclidean-product-borel}
$\mathscr{B}_{\mathbb{R}^n}=\bigotimes_1^n\mathscr{B}_{\mathbb{R}}$。
\end{corollary}
% SUPPLEMENT: id=v14-ch01-euclidean-borel-proof; parent=ch01-euclidean-product-borel; kind=proof+motivation+hypotheses; origin=independent-derivation
\begin{proof}
\textbf{补写证明（重构新增）。}实线在通常距离下由可数集 $\mathbb{Q}$ 稠密，故可分。对有限正整数 $n$，把上面的命题应用于 $n$ 个实线因子，得到相应乘积度量的 Borel 集族等于 $\bigotimes_1^n\mathscr{B}_{\mathbb{R}}$。通常 Euclid 距离与最大值距离满足
\begin{align*}
\max_{1\leqslant j\leqslant n}|x_j-y_j|
\leqslant |x-y|
\leqslant \sqrt{n}\max_{1\leqslant j\leqslant n}|x_j-y_j|.
\end{align*}
所以两种距离的开球互相包含适当缩放的球，产生同一开集族和同一 Borel $\sigma$-代数，结论也适用于通常的 $\mathbb{R}^n$。
\end{proof}
\begin{remark}
\textbf{动机与条件说明（重构新增）。}此推论把多维 Borel 可测性与逐坐标的实线 Borel 结构对应起来。$n$ 为有限正整数，使前述有限乘积命题及距离比较中的常数 $\sqrt n$ 适用；实线可分性由有理数稠密给出。这里比较的是 $\sigma$-代数，没有选取测度，也没有断言乘积测度或完备化后的集合族相等，这些是后续不同的问题。
\end{remark}
% END SUPPLEMENT: id=v14-ch01-euclidean-borel-proof


% SOURCE p=38; printed=23; object=ch01-elementary-family; region=lower
本节最后给出一个以后需要的技术性结果。
\begin{definition}[初等族]\label{definition:ch01-elementary-family}
称 $X$ 的子集族 $\mathscr{E}$ 为\textbf{初等族}，是指它满足：
\begin{itemize}
\item $\varnothing\in\mathscr{E}$；
\item 若 $E,F\in\mathscr{E}$，则 $E\cap F\in\mathscr{E}$；
\item 若 $E\in\mathscr{E}$，则 $E^{\mathrm{c}}$ 是 $\mathscr{E}$ 中有限个成员的不交并。
\end{itemize}
\end{definition}
% SUPPLEMENT: id=v14-ch01-elementary-family-remark; parent=ch01-elementary-family; kind=motivation+hypotheses; origin=independent-derivation
\begin{remark}
\textbf{动机与条件说明（重构新增）。}初等族保留容易描述的基本块，让有限不交拼接之后成为可做代数运算的集合族。包含空集保证退化部分也属于基本块；交封闭保证多个切分同时进行时，重叠部分仍是基本块；补集可写成有限不交基本块，使删除操作可以转化为有限分块。这三个条件针对同一个全集 $X$，补集尤其依赖这个环境。初等族本身不必对并封闭，也不必包含 $X$，但对 $\varnothing$ 应用补集条件，已保证 $X$ 能被有限个基本块拼成。有限与不交的要求分别控制块的数量及避免重复，正用于下一命题中的细分构造。
\end{remark}
% END SUPPLEMENT: id=v14-ch01-elementary-family-remark



% SOURCE p=38; printed=23; object=ch01-elementary-algebra; number=1.7; region=bottom
\begin{proposition}\label{proposition:ch01-elementary-algebra}
若 $\mathscr{E}$ 为初等族，则由 $\mathscr{E}$ 中成员的有限不交并构成的族 $\mathscr{A}$ 是代数。
\end{proposition}

% SOURCE p=39; printed=24; object=ch01-proof-elementary-algebra; region=top
\begin{proof}
若 $A,B\in\mathscr{E}$，且 $B^{\mathrm{c}}=\bigcup_1^JC_j$（$C_j\in\mathscr{E}$ 且两两不交），则 $A\setminus B=\bigcup_1^J(A\cap C_j)$，$A\cup B=(A\setminus B)\cup B$；这些并都是不交并，所以 $A\setminus B\in\mathscr{A}$ 且 $A\cup B\in\mathscr{A}$。由归纳法可得，若 $A_1,\ldots,A_n\in\mathscr{E}$，则 $\bigcup_1^nA_j\in\mathscr{A}$；事实上，由归纳假设可设 $A_1,\ldots,A_{n-1}$ 两两不交，于是 $\bigcup_1^nA_j=A_n\cup\bigcup_1^{n-1}(A_j\setminus A_n)$，这是不交并。为证明 $\mathscr{A}$ 对补运算封闭，设 $A_1,\ldots,A_n\in\mathscr{E}$，且 $A_m^{\mathrm{c}}=\bigcup_{j=1}^{J_m}B_m^j$，其中 $B_m^1,\ldots,B_m^{J_m}$ 是 $\mathscr{E}$ 中两两不交的成员。则
\begin{align*}
\Bigl(\bigcup_{m=1}^nA_m\Bigr)^{\mathrm{c}}
&=\bigcap_{m=1}^n\Bigl(\bigcup_{j=1}^{J_m}B_m^j\Bigr)\\
&=\bigcup\{B_1^{j_1}\cap\cdots\cap B_n^{j_n}:1\leqslant j_m\leqslant J_m,\ 1\leqslant m\leqslant n\},
\end{align*}
它属于 $\mathscr{A}$。
\end{proof}
% SUPPLEMENT: id=v14-ch01-elementary-algebra-proof; parent=ch01-elementary-algebra; kind=proof+motivation+hypotheses; origin=independent-derivation
\begin{proof}
\textbf{归纳及不交性细节的补全（重构新增）。}首先，任意有限个初等族成员的并都可以表示为有限不交基本块的并。具体给定 $E_1,\ldots,E_n\in\mathscr{E}$，置 $D_1=E_1$，并对 $k>1$ 置
\begin{align*}
D_k=E_k\cap\bigcap_{j=1}^{k-1}E_j^{\mathrm c}.
\end{align*}
每个补集 $E_j^{\mathrm c}$ 都是有限个两两不交的 $\mathscr{E}$ 成员之并。将有限交对这些有限并分配，$D_k$ 就成为有限个集合的并；每个集合是 $E_k$ 与有限个基本块的交，交封闭保证它仍在 $\mathscr{E}$ 中。同一个 $D_k$ 的两块若选择的某一补集分块指标不同，则它们在该坐标的两个不交基本块中，故交为空；因此所得并确实不交。不同 $D_k$ 之间也不交，因为后者删除了先前的 $E_j$。逐点按第一次出现的指标分配，给出 $\bigcup_kD_k=\bigcup_kE_k$。若某个补集为空，可用单个空基本块表示，避免空分配式的歧义。

现在取 $A,B\in\mathscr{A}$。它们各由有限个基本块组成，把这些基本块合在一个有限列表后应用刚证结论，得到 $A\cup B\in\mathscr{A}$；重复二元并即可得到有限并封闭。再设 $A=\bigcup_{m=1}^nA_m$ 为有限不交基本块表示，按照原文将各 $A_m^{\mathrm c}$ 分为有限不交的 $B_m^j$。分配式中的每个交块属于 $\mathscr{E}$；两组不同多重指标至少在某个 $m$ 上不同，相应两个 $B_m^j$ 不交，所以这些交块两两不交。故原文的补集公式确实给出 $A^{\mathrm c}\in\mathscr{A}$，而不只是一般有限并表示。空集本身属于 $\mathscr{E}\subseteq\mathscr{A}$，因此 $\mathscr{A}$ 非空。非空、有限并封闭、补封闭全部成立，$\mathscr{A}$ 是代数。
\end{proof}
\begin{remark}
\textbf{动机与条件说明（重构新增）。}由简单基本块构造代数时，不能只证明两个基本块的并可分解，还必须处理任意两个已拼接集合的并以及它们的补。补写构造用有限切分把这些情况统一起来：交封闭保证切出的块仍合法，补的有限分块保证产生的块总数有限，分块不交性保证最终表示符合 $\mathscr{A}$ 的定义。这里不要求可数拼接，故结论是代数；没有据此断言 $\sigma$-代数封闭性。
\end{remark}
% END SUPPLEMENT: id=v14-ch01-elementary-algebra-proof


% SOURCE p=39; printed=24; object=ch01-sec02-exercises; region=middle
\Needspace{14\baselineskip}
\subsection*{习题}

% SOURCE p=39; printed=24; object=ch01-exercise-rings; number=1; region=middle
\begin{exercise}\label{exercise:ch01-rings}

\begin{definition*}[环与 $\sigma$-环]
若集合族 $\mathscr{R}\subseteq\mathscr{P}(X)$ 对有限并和差运算封闭（即若 $E_1,\ldots,E_n\in\mathscr{R}$，则 $\bigcup_1^nE_j\in\mathscr{R}$；若 $E,F\in\mathscr{R}$，则 $E\setminus F\in\mathscr{R}$），则称 $\mathscr{R}$ 为\textbf{环}。对可数并封闭的环称为 \textbf{$\sigma$-环}。
\end{definition*}
% SUPPLEMENT: id=v14-ch01-rings-definition-remark; parent=ch01-exercise-rings; kind=motivation+hypotheses+source-qualification; origin=independent-derivation
\begin{remark}
\textbf{局部定义的动机与条件说明（重构新增）。}环允许做有限集合拼接与相对删除，却不要求整个环境 $X$ 可用；这适合先在局部集合上工作。差运算只需要已有的两个成员，所以比相对于固定 $X$ 取补更弱；包含 $X$ 后，两者由 $E^{\mathrm c}=X\setminus E$ 联系起来。可数并的要求则使 $\sigma$-环能处理无限次拼接。以上各集合都位于同一 $X$ 中，避免混用不同环境。

\textbf{原文约定说明（重构新增）。}原题定义未明写环非空。若“有限并”仅检验正整数个成员，空族 $\mathscr{R}=\varnothing$ 也会空真地满足所列封闭条件，却使后面的\hyperref[item:ch01-rings-complements]{补集扩充结论}变为空族是 $\sigma$-代数，结论不成立。若“有限并”包括空并，则 $\varnothing\in\mathscr{R}$ 已排除这个问题。以下采用通常的非空环约定；等价地，可在所列封闭条件上明确补加 $\mathscr{R}\ne\varnothing$。非空时取 $E\in\mathscr{R}$，差封闭立即给出 $\varnothing=E\setminus E\in\mathscr{R}$。

本题交的封闭性针对非空指标族；不能要求环包含空族在 $X$ 中的交，因为该交为 $X$，而环未必含 $X$。例如 $X=\mathbb{N}$ 上所有有限子集组成非空环，却不含 $X$。若保留允许空环的解释，下面其余结论仍可按原量词理解；空环的补集扩充结论则必须按上述反例单独排除。
\end{remark}
% END SUPPLEMENT: id=v14-ch01-rings-definition-remark


\begin{enumerate}[label=\alph*.,ref=(\alph*)]
\item\label{item:ch01-rings-intersections} 环（分别地，$\sigma$-环）对有限（分别地，可数）交封闭。
\item\label{item:ch01-rings-algebras} 若 $\mathscr{R}$ 为环（分别地，$\sigma$-环），则 $\mathscr{R}$ 为代数（分别地，$\sigma$-代数）当且仅当 $X\in\mathscr{R}$。
\item\label{item:ch01-rings-complements} 若 $\mathscr{R}$ 为 $\sigma$-环，则 $\{E\subseteq X:E\in\mathscr{R}\text{ 或 }E^{\mathrm{c}}\in\mathscr{R}\}$ 是 $\sigma$-代数。
\item\label{item:ch01-rings-localization} 若 $\mathscr{R}$ 为 $\sigma$-环，则 $\{E\subseteq X:\text{对所有 }F\in\mathscr{R},\ E\cap F\in\mathscr{R}\}$ 是 $\sigma$-代数。
\end{enumerate}
\end{exercise}
% SUPPLEMENT: id=v14-ch01-rings-solutions; parent=ch01-exercise-rings; kind=solution; origin=independent-derivation
\begin{solution}
\phantomsection\label{solution:ch01-rings}
\textbf{习题解答（重构新增）。}按上方已明确的非空环约定证明各问；全部补集均相对于 $X$。
\begin{enumerate}[label=\alph*.,ref=(\alph*)]
\item\label{item:ch01-rings-solution-intersections} 若 $E,F\in\mathscr{R}$，则
\begin{align*}
E\cap F=E\setminus(E\setminus F)\in\mathscr{R}.
\end{align*}
对非空有限族反复使用这个二元等式即得交封闭。若 $\mathscr{R}$ 是 $\sigma$-环，且 $E_n\in\mathscr{R}$ 对所有 $n\geqslant1$ 成立，则每个 $E_1\setminus E_n$ 在环中，其可数并也在环中，所以
\begin{align*}
\bigcap_{n=1}^{\infty}E_n
=E_1\setminus\bigcup_{n=1}^{\infty}(E_1\setminus E_n)
\in\mathscr{R}.
\end{align*}
等式逐点成立：一个点在右端，恰表示它属于 $E_1$ 且没有任何一个 $E_n$ 把它排除。有限指标族已处理，可数无限指标族经枚举得到此式。
\item\label{item:ch01-rings-solution-algebras} 若 $\mathscr{R}$ 为代数，则由非空性可取 $E\in\mathscr{R}$；补与有限并封闭给出 $X=E\cup E^{\mathrm c}\in\mathscr{R}$。反之，若 $X\in\mathscr{R}$，差封闭给出对每个 $E\in\mathscr{R}$ 有 $E^{\mathrm c}=X\setminus E\in\mathscr{R}$；连同已有有限并封闭和非空性，这正是代数定义。若原环还是 $\sigma$-环，则同一论证加上已有可数并封闭给出 $\sigma$-代数；必要方向因每个 $\sigma$-代数也是代数而同样成立。
\item\label{item:ch01-rings-solution-complements} 记待证集合族为 $\mathscr{S}$。由 $\varnothing\in\mathscr{R}$ 可知 $X,\varnothing\in\mathscr{S}$。若 $E\in\mathscr{S}$，则 $E\in\mathscr{R}$ 或 $E^{\mathrm c}\in\mathscr{R}$；交换 $E,E^{\mathrm c}$ 并用 $(E^{\mathrm c})^{\mathrm c}=E$，得到 $E^{\mathrm c}\in\mathscr{S}$。

设 $E_n\in\mathscr{S}$，令 $G=\bigcup_nE_n$。若所有 $E_n$ 都在 $\mathscr{R}$ 中，则 $G\in\mathscr{R}\subseteq\mathscr{S}$。否则存在 $k$ 使 $E_k\notin\mathscr{R}$，于是 $C=E_k^{\mathrm c}\in\mathscr{R}$。对每个 $n$，若 $E_n\in\mathscr{R}$，则 $C\cap E_n^{\mathrm c}=C\setminus E_n\in\mathscr{R}$；若 $E_n\notin\mathscr{R}$，则 $E_n^{\mathrm c}\in\mathscr{R}$，由已证交封闭也得 $C\cap E_n^{\mathrm c}\in\mathscr{R}$。因此由本解答的\hyperref[item:ch01-rings-solution-intersections]{可数交结论}，
\begin{align*}
G^{\mathrm c}=\bigcap_{n=1}^{\infty}(C\cap E_n^{\mathrm c})\in\mathscr{R}.
\end{align*}
这里左侧已包含于 $C=E_k^{\mathrm c}$，所以加上 $C$ 不改变交。故 $G\in\mathscr{S}$。非空、补封闭、可数并封闭全部成立，$\mathscr{S}$ 为 $\sigma$-代数。若允许 $\mathscr{R}=\varnothing$，则本问按原字面失败，反例已在上方说明，不能给出相反结论的证明。
\item\label{item:ch01-rings-solution-localization} 记待证集合族为 $\mathscr{L}$。对任意 $F\in\mathscr{R}$，$X\cap F=F\in\mathscr{R}$，故 $X\in\mathscr{L}$。若 $E\in\mathscr{L}$，则对每个 $F\in\mathscr{R}$，
\begin{align*}
E^{\mathrm c}\cap F=F\setminus(E\cap F)\in\mathscr{R},
\end{align*}
因为 $E\cap F\in\mathscr{R}$ 且环对差封闭。因此 $E^{\mathrm c}\in\mathscr{L}$。若 $E_n\in\mathscr{L}$，则对每个 $F\in\mathscr{R}$，
\begin{align*}
\left(\bigcup_{n=1}^{\infty}E_n\right)\cap F
=\bigcup_{n=1}^{\infty}(E_n\cap F)\in\mathscr{R},
\end{align*}
这里使用 $\sigma$-环的可数并封闭。所以 $\bigcup_nE_n\in\mathscr{L}$，得到 $\sigma$-代数。若 $\mathscr{R}=\varnothing$，检验条件对每个 $E\subseteq X$ 都空真，此时 $\mathscr{L}=\mathscr{P}(X)$，本问仍成立。
\end{enumerate}
\end{solution}
% END SUPPLEMENT: id=v14-ch01-rings-solutions


% SOURCE p=39; printed=24; object=ch01-exercise-borel-generators-proof; number=2; region=middle
\begin{exercise}\label{exercise:ch01-borel-generators-proof}
补全\cref{proposition:ch01-real-borel-generators}的证明。
\end{exercise}
% SUPPLEMENT: id=v14-ch01-borel-generators-solution; parent=ch01-exercise-borel-generators-proof; kind=solution; origin=independent-derivation
\begin{solution}
\phantomsection\label{solution:ch01-borel-generators-proof}
\textbf{习题解答（重构新增）。}沿用\cref{proposition:ch01-real-borel-generators}中的八个集合族。首先，各生成元都是 Borel 集：开区间、开射线是开集，闭区间、闭射线是闭集；两种半开区间满足
\begin{align*}
(a,b]&=\bigcap_{n=1}^{\infty}(a,b+n^{-1}),&
[a,b)&=\bigcap_{n=1}^{\infty}(a-n^{-1},b).
\end{align*}
这两个等式逐点排除左端点或右端点外的点，同时保留所需闭端点。因此每个 $\mathscr{M}(\mathscr{E}_j)$ 都包含于 $\mathscr{B}_{\mathbb{R}}$。

反向只需从各族生成任意有界开区间 $(a,b)$：实直线上的每个开集可写为有界开区间的可数并，具体可取全部包含于该开集的有理端点开区间；这样的区间可数，而且由开性及有理数稠密性覆盖每个点。下面给出八种生成的全部验证。
\begin{enumerate}[label=\alph*.,ref=(\alph*)]
\item 对 $\mathscr{E}_1$，所需开区间本来就是生成元。
\item 对 $\mathscr{E}_2$，令 $\delta_n=(b-a)/(n+3)$。则 $0<\delta_n<(b-a)/2$ 且 $\delta_n\to0$，故
\begin{align*}
(a,b)=\bigcup_{n=1}^{\infty}[a+\delta_n,b-\delta_n].
\end{align*}
每个闭区间的左端严格小于右端，确实属于题设的生成族；并集包含于 $(a,b)$。反之，$a<x<b$ 时，可取 $n$ 使 $\delta_n<\min(x-a,b-x)$，则 $x$ 落入该闭区间。这样避免使用端点倒置或退化的区间。
\item 对 $\mathscr{E}_3$ 和 $\mathscr{E}_4$，使用同一 $\delta_n$，有
\begin{align*}
(a,b)&=\bigcup_{n=1}^{\infty}(a,b-\delta_n],&
(a,b)&=\bigcup_{n=1}^{\infty}[a+\delta_n,b).
\end{align*}
每个成员均为合法非退化生成元，且包含于 $(a,b)$。给定 $x\in(a,b)$，第一式取 $\delta_n<b-x$，第二式取 $\delta_n<x-a$，分别使 $x$ 属于右端某一成员，证明等式。
\item 对 $\mathscr{E}_5$，先由
\begin{align*}
[b,\infty)&=\bigcap_{n=1}^{\infty}(b-n^{-1},\infty),&
(a,b)&=(a,\infty)\setminus[b,\infty)
\end{align*}
得到开区间。交式包含所有 $x\geqslant b$；若 $x<b$，选 $n$ 使 $n^{-1}<b-x$ 即排除 $x$。对 $\mathscr{E}_6$，对应地使用
\begin{align*}
(-\infty,a]&=\bigcap_{n=1}^{\infty}(-\infty,a+n^{-1}),&
(a,b)&=(-\infty,b)\setminus(-\infty,a].
\end{align*}
交式包含所有 $x\leqslant a$；若 $x>a$，选 $n^{-1}<x-a$ 即排除它。差运算在 $\sigma$-代数中由补和交实现。
\item 对 $\mathscr{E}_7$，使用
\begin{align*}
(a,\infty)&=\bigcup_{n=1}^{\infty}[a+n^{-1},\infty),&
(a,b)&=(a,\infty)\setminus[b,\infty).
\end{align*}
并式每项都位于 $(a,\infty)$；若 $x>a$，选 $n^{-1}<x-a$ 则该项包含 $x$。对 $\mathscr{E}_8$，使用
\begin{align*}
(-\infty,b)&=\bigcup_{n=1}^{\infty}(-\infty,b-n^{-1}],&
(a,b)&=(-\infty,b)\setminus(-\infty,a].
\end{align*}
若 $x<b$，取 $n^{-1}<b-x$ 即可，反向包含直接由端点得出。
\end{enumerate}
所以每一族生成的 $\sigma$-代数都含所有开集，按 Borel $\sigma$-代数的最小性包含 $\mathscr{B}_{\mathbb{R}}$。结合前面的反向包含，八种情形全部得到相等。
\end{solution}
% END SUPPLEMENT: id=v14-ch01-borel-generators-solution


% SOURCE p=39; printed=24; object=ch01-exercise-infinite-sigma-algebra; number=3; region=middle
\begin{exercise}\label{exercise:ch01-infinite-sigma-algebra}
设 $\mathscr{M}$ 为无限 $\sigma$-代数。
\begin{enumerate}[label=\alph*.,ref=(\alph*)]
\item\label{item:ch01-infinite-sigma-disjoint} $\mathscr{M}$ 包含一个由两两不交集合组成的无限序列。
\item\label{item:ch01-infinite-sigma-cardinality} $\card(\mathscr{M})\geqslant\mathfrak{c}$。
\end{enumerate}
\end{exercise}
% SUPPLEMENT: id=v14-ch01-infinite-sigma-solution; parent=ch01-exercise-infinite-sigma-algebra; kind=solution; origin=independent-derivation
\begin{solution}
\phantomsection\label{solution:ch01-infinite-sigma-algebra}
\textbf{习题解答（重构新增）。}
\begin{enumerate}[label=\alph*.,ref=(\alph*)]
\item\label{item:ch01-infinite-sigma-solution-disjoint} 实际构造两两不交且每项非空的序列，这比允许重复空集的字面结论更强，也是下一问所需要的形式。对 $B\in\mathscr{M}$，记
\begin{align*}
\mathscr{M}_B=\{E\cap B:E\in\mathscr{M}\}.
\end{align*}
这是用于计数的辅助集合族，其每个成员也在 $\mathscr{M}$ 中。若 $\mathscr{M}_B$ 无限，则其中存在满足 $\varnothing\ne C\subsetneq B$ 的成员，否则它至多含 $\varnothing,B$ 两个成员。置 $D=B\setminus C$；$C,D$ 都非空且都在 $\mathscr{M}$ 中。映射
\begin{align*}
H\longmapsto(H\cap C,H\cap D),\qquad H\in\mathscr{M}_B,
\end{align*}
为从 $\mathscr{M}_B$ 到 $\mathscr{M}_C\times\mathscr{M}_D$ 的单射，因为 $H=(H\cap C)\cup(H\cap D)$。故 $\mathscr{M}_C,\mathscr{M}_D$ 不可能同时有限，否则有限集的乘积有限，迫使 $\mathscr{M}_B$ 有限。

从 $B_0=X$ 开始，此时 $\mathscr{M}_{B_0}=\mathscr{M}$ 无限。按上面的步骤，把 $B_{n-1}$ 分为两个非空可测部分，选择其中仍有无限辅助集合族的一部分作 $B_n$，把另一部分记为 $A_n$。这种可数递归选择使用本书已采用的选择公理。于是
\begin{align*}
B_n\subseteq B_{n-1},\qquad
A_n=B_{n-1}\setminus B_n\in\mathscr{M},\qquad A_n\ne\varnothing.
\end{align*}
若 $m>n$，则 $A_m\subseteq B_{m-1}\subseteq B_n$，而 $A_n\cap B_n=\varnothing$，所以 $A_n\cap A_m=\varnothing$。得到所需无限序列。
\item\label{item:ch01-infinite-sigma-solution-cardinality} 取刚构造的序列。对 $S\subseteq\mathbb{N}$，令
\begin{align*}
\Phi(S)=\bigcup_{n\in S}A_n.
\end{align*}
这是可数并（$S=\varnothing$ 时为空集），故属于 $\mathscr{M}$。若 $S\ne T$，可交换二者使存在 $k\in S\setminus T$。因 $A_k$ 非空，取 $x\in A_k$；两两不交性给出 $x\in\Phi(S)$ 而 $x\notin\Phi(T)$。所以 $\Phi$ 为单射。

为直接证明所需连续统下界，枚举有理数为 $\{q_n\}_{n\geqslant1}$，其中每个有理数都出现。这样的枚举可由 $\mathbb{Q}=\{p/m:p\in\mathbb{Z},\ m\in\mathbb{N}\}$ 及整数对的可数性得到。映射
\begin{align*}
\Psi:\mathbb{R}\longrightarrow\mathscr{P}(\mathbb{N}),\qquad
\Psi(t)=\{n:q_n<t\}
\end{align*}
是单射：若 $s<t$，有理数稠密性给出 $s<q_k<t$，故 $k\in\Psi(t)\setminus\Psi(s)$。因此 $\Phi\circ\Psi$ 是 $\mathbb{R}$ 到 $\mathscr{M}$ 的单射，按基数比较的定义，$\card(\mathscr{M})\geqslant\mathfrak{c}$。这个证明不需要先证明二进制展开的唯一性。
\end{enumerate}
\end{solution}
% END SUPPLEMENT: id=v14-ch01-infinite-sigma-solution


% SOURCE p=39; printed=24; object=ch01-exercise-increasing-unions; number=4; region=middle-lower
\begin{exercise}\label{exercise:ch01-increasing-unions}
代数 $\mathscr{A}$ 为 $\sigma$-代数，当且仅当 $\mathscr{A}$ 对可数递增并封闭（即若 $\{E_j\}_1^\infty\subseteq\mathscr{A}$ 且 $E_1\subseteq E_2\subseteq\cdots$，则 $\bigcup_1^\infty E_j\in\mathscr{A}$）。
\end{exercise}
% SUPPLEMENT: id=v14-ch01-increasing-unions-solution; parent=ch01-exercise-increasing-unions; kind=solution; origin=independent-derivation
\begin{solution}
\phantomsection\label{solution:ch01-increasing-unions}
\textbf{习题解答（重构新增）。}若 $\mathscr{A}$ 是 $\sigma$-代数，则它对任意可数并封闭，当然也对其中的递增并封闭。

反之，假设代数 $\mathscr{A}$ 对可数递增并封闭。任取 $E_n\in\mathscr{A}$，令 $F_n=\bigcup_{j=1}^nE_j$。有限并封闭保证每个 $F_n\in\mathscr{A}$，且 $F_n\subseteq F_{n+1}$，所以 $F=\bigcup_nF_n\in\mathscr{A}$。若 $x\in F$，则对某个 $n$、某个 $j\leqslant n$ 有 $x\in E_j$，故 $x\in\bigcup_jE_j$；反之，若 $x\in E_j$，则 $x\in F_j\subseteq F$。因而 $F=\bigcup_jE_j$。这证明任意可数并都在 $\mathscr{A}$ 中，加上原有代数性质，$\mathscr{A}$ 是 $\sigma$-代数。有限并假设正用于把任意序列改造成递增序列，未对原来的 $E_n$ 暗加单调性。
\end{solution}
% END SUPPLEMENT: id=v14-ch01-increasing-unions-solution


% SOURCE p=39; printed=24; object=ch01-exercise-countable-generators; number=5; region=lower
\begin{exercise}\label{exercise:ch01-countable-generators}
若 $\mathscr{M}$ 为由 $\mathscr{E}$ 生成的 $\sigma$-代数，则 $\mathscr{M}$ 是当 $\mathscr{F}$ 遍历 $\mathscr{E}$ 的所有可数子集时，由 $\mathscr{F}$ 生成的各 $\sigma$-代数的并。（提示：证明后一对象是 $\sigma$-代数。）
\end{exercise}
% SUPPLEMENT: id=v14-ch01-countable-generators-solution; parent=ch01-exercise-countable-generators; kind=solution; origin=independent-derivation
\begin{solution}
\phantomsection\label{solution:ch01-countable-generators}
\textbf{习题解答（重构新增）。}记
\begin{align*}
\mathscr{U}=\bigcup\{\mathscr{M}(\mathscr{F}):
\mathscr{F}\subseteq\mathscr{E},\ \mathscr{F}\text{ 可数}\}.
\end{align*}
因为空集是 $\mathscr{E}$ 的可数子集，$\mathscr{M}(\varnothing)\subseteq\mathscr{U}$，所以 $X,\varnothing\in\mathscr{U}$，即使 $\mathscr{E}$ 本身为空也成立。

若 $E\in\mathscr{U}$，则存在可数 $\mathscr{F}\subseteq\mathscr{E}$ 使 $E\in\mathscr{M}(\mathscr{F})$。后者对补封闭，故 $E^{\mathrm c}\in\mathscr{M}(\mathscr{F})\subseteq\mathscr{U}$。若 $E_n\in\mathscr{U}$，对每个 $n$ 选取可数 $\mathscr{F}_n\subseteq\mathscr{E}$，使 $E_n\in\mathscr{M}(\mathscr{F}_n)$；这里允许使用本书的选择公理。可数个可数集的并
\begin{align*}
\mathscr{F}=\bigcup_{n=1}^{\infty}\mathscr{F}_n
\end{align*}
仍可数，且包含于 $\mathscr{E}$。由 $\mathscr{F}_n\subseteq\mathscr{F}$ 和生成族的最小性，每个 $\mathscr{M}(\mathscr{F}_n)\subseteq\mathscr{M}(\mathscr{F})$，所以所有 $E_n$ 均在同一个 $\sigma$-代数 $\mathscr{M}(\mathscr{F})$ 中。于是 $\bigcup_nE_n\in\mathscr{M}(\mathscr{F})\subseteq\mathscr{U}$。这证明 $\mathscr{U}$ 是 $\sigma$-代数；不能把一般 $\sigma$-代数之并误认为自动具有这个性质，关键在于能合并可数个生成族。

对每个 $E\in\mathscr{E}$，单元素族 $\{E\}$ 可数且 $E\in\mathscr{M}(\{E\})\subseteq\mathscr{U}$，故 $\mathscr{E}\subseteq\mathscr{U}$。由最小性，$\mathscr{M}(\mathscr{E})\subseteq\mathscr{U}$。另一方面，每个参与取并的 $\mathscr{F}$ 都包含于 $\mathscr{E}$，所以 $\mathscr{M}(\mathscr{F})\subseteq\mathscr{M}(\mathscr{E})$，取并得 $\mathscr{U}\subseteq\mathscr{M}(\mathscr{E})$。两向包含给出结论。
\end{solution}
% END SUPPLEMENT: id=v14-ch01-countable-generators-solution

\end{document}
